\scanpage{23}
Les éléments fixes sous l'involution $u\mapsto u^\varphi$ correspondent
alors aux formes \unsure{symm.|sym.} resp. alternées
(suivant $p$ pair ou impair). De plus, on a
\[
\left\{\begin{aligned}
(u+v)^\varphi&=u^\varphi+v^\varphi\\
(\lambda u)^\varphi&=\lambda u^\varphi && \lambda\in Q\\
(u^\varphi)^\varphi&=u\\
(uv)^\varphi&=\deleted{\unsure{\varphi}}v^\varphi u^\varphi
\end{aligned}\right.
\]

Si on remplace $\varphi$ par une forme $\psi$ de $=$
type (pair resp. impair) on aura $\psi\longleftrightarrow a$
\[
\psi(x,y)\deleted{\unsure{(x,\psi)}}=\varphi(ax,y)
\]
et l'hyp. sur $\psi$ s'exprime, comme on a dit
par $a^\varphi=a$. On aura de plus
\[
u^\psi=\deleted{\unsure{a^\varphi}}a^{-1}u^\varphi a\qquad\text{($a$ étant inv.).}
\]
donc l'\uline{involution de} $A=\mathrm{End}(M)$ \uline{est
bien déterminée mod autom intérieur}
(indépendamt du choix de $\varphi$, ou $\psi$), une
fois si $\varphi$ est \unsure{fixée|fixé}, mod aut. int.
par un élément \deleted{de $A$ symétrique} \inserted{$\varphi$-hermitien} de $A$.
% archive sheet number: 22
